\documentclass[10pt]{article}
\usepackage[T1]{fontenc}
\usepackage[utf8]{inputenc}
\usepackage[french]{babel}
\usepackage{amsmath,amssymb,mathrsfs,mathpazo}
\usepackage{geometry}
\geometry{a4paper,margin=2cm}
\usepackage{graphicx,array,tabularx,tabularray,colortbl,multirow,multicol,booktabs,enumitem,pifont,ulem,fancybox,diagbox,cancel}
\usepackage{pstricks,pstricks-add,pst-plot,pst-tree,pst-node,pst-text,pst-func,pst-math,pst-bspline,pst-3dplot,pst-eucl,pst-solides3d}
\let\pscalculate\undefined
\usepackage{pst-bezier}
\usepackage{tikz,pgfplots,tkz-tab}
\usetikzlibrary{arrows,arrows.meta,patterns,calc,positioning,trees,decorations.markings,shapes.geometric,shapes.misc,shapes.arrows,intersections,decorations.pathmorphing,decorations.pathreplacing,backgrounds,shadows,babel}
\pgfplotsset{compat=1.18}
\usepgfplotslibrary{fillbetween}
\usepackage[np]{numprint}
\usepackage{listings,makecell,pas-tableur,twoopt}
\let\barre\undefined
\usepackage[pstricks]{tablvar}
\definecolor{darkgreen}{RGB}{0,100,0}
\providecommand{\red}{\color{red}}
\providecommand{\blue}{\color{blue}}
\providecommand{\green}{\color{green}}
\providecommand{\black}{\color{black}}
\providecommand{\white}{\color{white}}
\providecommand{\gray}{\color{gray}}
\providecommand{\cyan}{\color{cyan}}
\providecommand{\magenta}{\color{magenta}}
\providecommand{\yellow}{\color{yellow}}
\providecommand{\violet}{\color{violet}}
\providecommand{\orange}{\color{orange}}
\providecommand{\euro}{\texteuro}
\providecommand{\EUR}{\texteuro}
\providecommand{\ieme}{\textsuperscript{e}}
\providecommand{\degres}{\ensuremath{^\circ}}
\providecommand{\degree}{\ensuremath{^\circ}}
\providecommand{\vv}[1]{\overrightarrow{#1}}
\providecommand{\Coord}[1]{#1}
\providecommand{\decosix}{}
\providecommand{\starredbullet}{\ensuremath{\bullet}}
\pagestyle{empty}
\setlength{\parindent}{0pt}
\usepackage{adjustbox}\DeclareUnicodeCharacter{2236}{:}\DeclareUnicodeCharacter{2190}{\ensuremath{\leftarrow}}
\providecommand{\R}{\mathbb{R}}
\providecommand{\N}{\mathbb{N}}
\providecommand{\Z}{\mathbb{Z}}
\providecommand{\Q}{\mathbb{Q}}
\providecommand{\D}{\mathbb{D}}
\providecommand{\C}{\mathbb{C}}
\providecommand{\e}{{\mathrm{e}}}
\providecommand{\vect}[1]{\overrightarrow{#1}}
\providecommand{\vectt}[1]{\overrightarrow{\mathrm{#1}}}
\providecommand{\vvt}[1]{\overrightarrow{\mathrm{#1}}}
\providecommand{\np}[2][]{{#2~#1}}
\providecommand{\vv}[1]{{\overrightarrow{#1}}}
\renewcommand{\d}{{\mathrm{d}}}
\providecommand{\degres}{{^\circ}}
\providecommand{\degre}{{^\circ}}
\providecommand{\blue}{\color{blue}}
\providecommand{\red}{\color{red}}
\providecommand{\enskip}{\quad}
\providecommand{\llbracket}{\lbrack\!\lbrack}
\providecommand{\rrbracket}{\rbrack\!\rbrack}
\providecommand{\euro}{{\text{€}}}
\providecommand{\no}{{\text{n°}}}
\providecommand{\interval}[2]{{[#1;#2]}}
\providecommand{\intervFO}[2]{{[#1;#2[}}
\providecommand{\dotfill}{\ldots}

\providecommand{\ds}{\displaystyle}
\providecommand{\textding}[1]{#1}
\providecommand{\Oij}{\(\mathrm{O};\vec{\imath},\vec{\jmath}\)}
\providecommand{\Oijk}{\(\mathrm{O};\vec{\imath},\vec{\jmath},\vec{k}\)}
\providecommand{\Ouv}{\(\mathrm{O};\vec{u},\vec{v}\)}
\providecommand{\intervFF}[2]{\ensuremath{[#1;#2]}}
\renewcommand{\euro}{\ifmmode\text{\texteuro}\else\texteuro\fi}
\DeclareUnicodeCharacter{00B0}{\ifmmode^{\circ}\else\textdegree\fi}
\DeclareUnicodeCharacter{20AC}{\ifmmode\text{\texteuro}\else\texteuro\fi}
% Énoncés : Éducation nationale. Transcription et corrigés : APMEP, Denis Vergès et contributeurs. Découpage et classement : CoopMaths, Sébastien Lozano, Éric Elter et contributeurs. Adaptation de la présentation : Maths Explorer.
% CC BY-SA : https://coopmaths.fr/www/about/
% https://coopmaths.fr/alea/static/bac/2025/tex/bac_2025_06_sujet1_polynesie_1.tex
\begin{document}
\noindent\textit{APMEP / CoopMaths. CC BY-SA. Présentation adaptée par Maths Explorer.}\par\medskip
\medskip

Une équipe américaine a cartographié pour la première fois les allergies alimentaires chez l'enfant aux États-Unis en 2020. L'étude, publiée dans la revue \emph{Clinical Pediatries}, révèle une différence nette entre les zones rurales et les zones urbaines.

\smallskip

On sait qu'en 2020, 17\,\% de la population des États-Unis habite en zone rurale et 83\,\% en zone urbaine.

\smallskip

L'étude menée montre que parmi les enfants des États-Unis vivant en zone rurale, il y en a 6,2\,\% qui sont atteints d'allergie alimentaire.

L'étude révèle aussi que 9\,\% des enfants des États-Unis sont atteints d'allergie alimentaire. 

\smallskip

Pour un évènement $E$ quelconque, on note $p(E)$ sa probabilité et $\overline{E}$ son évènement contraire.

\textbf{Sauf mention contraire, les probabilités seront données sous forme exacte.}

\medskip

\textbf{Partie A}

\medskip

On interroge au hasard un enfant dans la population des États-Unis et on note :

\begin{itemize}[label=$\bullet~$]
\item $R$ L'évènement: \og l'enfant interrogé habite en zone rurale \fg ;
\item $A$ L'évènement: \og l'enfant interrogé est atteint d'allergie alimentaire \fg.
\end{itemize}

\medskip

\begin{enumerate}
\item Traduire cette situation à l'aide d'un arbre de probabilité. Cet arbre pourra être complété par la suite.
\item 
	\begin{enumerate}
		\item Calculer la probabilité que l'enfant interrogé habite en zone rurale et soit atteint d'allergie alimentaire,
		\item En déduire la probabilité que l'enfant interrogé habite en zone urbaine et soit atteint d'allergie alimentaire.
		\item L'enfant interrogé habite en zone urbaine. Quelle est la probabilité qu'il soit atteint d'allergie alimentaire ? Arrondir le résultat à $10^{-4}$.
	\end{enumerate}
\end{enumerate}

\medskip

\textbf{Partie B}

\medskip

On réalise une étude en interrogeant au hasard $100$ enfants des États-Unis.

On admet que ce choix se ramène à des tirages successifs indépendants avec remise.

On note $X$ la variable aléatoire donnant le nombre d'enfants atteints d'allergie alimentaire dans l'échantillon considéré.

\medskip

\begin{enumerate}
\item Justifier que la variable aléatoire $X$ suit une loi binomiale dont on précisera les paramètres.
\item Quelle est la probabilité qu'au moins 10 enfants parmi les $100$ interrogés soient atteints d'allergie alimentaire ? Arrondir le résultat à $10^{-4}$.
\end{enumerate}

\medskip

\textbf{Partie C}

\medskip

On s'intéresse à un échantillon de 20 enfants atteints d'allergie alimentaire choisis au
hasard.

\smallskip

L'âge d'apparition des premiers symptômes allergiques de ces 20 enfants est modélisé par les variables aléatoires $A_1,\: A_2,\ldots, ,A_{20}$. On admet que ces variables aléatoires sont indépendantes et suivent la même loi d'espérance $4$ et de variance $2,25$.
On considère la variable aléatoire:

\[M_{20} = \dfrac{A_1 + A_2 + \ldots + A_{20}}{20}.\]


\begin{enumerate}
\item Que représente la variable aléatoire $M_{20}$ dans le contexte de l'exercice ?
\item Déterminer l'espérance et la variance de $M_{20}$.
\item Justifier, à l'aide de l'inégalité de concentration, que

\[P(2 < M_{20} < 6) > 0,97.\]

Interpréter ce résultat dans le contexte de l'exercice.
\end{enumerate}

\bigskip

\end{document}
